\documentclass[11pt,a4paper]{article}

\usepackage{amsmath,amssymb,amsthm}
\usepackage{hyperref}
\usepackage{geometry}
\usepackage{doi}
\geometry{margin=2.5cm}
\setlength{\emergencystretch}{2em}

\newtheorem{theorem}{Theorem}[section]
\newtheorem{lemma}[theorem]{Lemma}
\newtheorem{corollary}[theorem]{Corollary}
\newtheorem{proposition}[theorem]{Proposition}
\newtheorem{definition}[theorem]{Definition}
\newtheorem{remark}[theorem]{Remark}

\newcommand{\bstar}{\beta^{*}}
\newcommand{\ImH}{\mathrm{Im}\,\mathbb{H}}
\newcommand{\su}{\mathfrak{su}}
\newcommand{\SU}{\mathrm{SU}}
\newcommand{\SO}{\mathrm{SO}}
\newcommand{\SL}{\mathrm{SL}}
\newcommand{\Heis}{\mathrm{Heis}_{3}(\mathbb{Z}/q\mathbb{Z})}
\newcommand{\sigc}{\Sigma_{c}}
\newcommand{\dpair}{\delta_{\mathrm{pair}}}

\title{Three Neutral Directions from a Spinor Carrier:\\
Conditional Status of the Threefold Admissibility Threshold\\
\large O23 --- Spectral Admissibility Sub-programme}

\author{J\'{e}r\^{o}me Beau\\
\small Independent Researcher}

\date{2026\\[2pt]\small Version 2.1}

\hypersetup{
  pdftitle={Three Neutral Directions from a Spinor Carrier: Conditional Status of the Threefold
    Admissibility Threshold},
  pdfauthor={J\'{e}r\^{o}me Beau},
  pdfsubject={Spectral Admissibility Sub-programme (O23), Cosmochrony},
  pdfkeywords={spectral admissibility; spinor carrier; quaternions; Born-Infeld constraint;
    SU(2) representations; imaginary quaternions; Weil representation; ADE classification;
    Borel subgroup; binary polyhedral groups; carrier selection}
}

\begin{document}

  \maketitle

  \begin{abstract}
    O21~\cite{Beau2026a25} defined the observable shell $n_{3}$ via the threshold condition
    $\Sigma_{c}(n_{3}) = 3$ and used it as structural input for the fibre-level transfer
    $\dpair \to \bstar$.
    The present paper determines the exact status of the value three in that threshold.
    The positive result is conditional: if the admissible neutral sector is carried by an
    irreducible two-dimensional $\SU(2)$-valued representation $V_{\rho} \cong \mathbb{C}^{2}$,
    then its traceless neutral sector is the real Lie algebra $\su(2) \cong \ImH$, of real
    dimension exactly three.
    This is a representation-theoretic theorem about the supplied carrier, and it
    explains why a spinorial neutral sector carries three independent directions.
    Two bridges required to convert this theorem into a derivation of
    $\Sigma_{c}(n_{3}) = 3$ are open and are stated here explicitly.
    On the first of them a second, unconditional theorem closes one candidate source exactly.
    The fingerprint vectors of the O12 shell-span filtration are pure Fourier modes, so every
    proper shell span is a coordinate subspace whose Weyl support is a single line, and its
    exact stabiliser in the Weil image of $\SL(2,\mathbb{Z}/q\mathbb{Z})$ is the metacyclic
    group $U \rtimes M_{n}$ inside the corresponding Borel subgroup.
    Because no binary polyhedral group embeds in a metacyclic Borel, that filtration selects
    none of $2T$, $2O$, $2I$, and therefore supplies no two-dimensional spinor carrier, for
    every odd prime, every generic block and every proper level --- including the primes
    $q \equiv \pm 1 \pmod 5$ where $2I$ is genuinely available.
    This closes one source of the carrier, not the carrier itself, and is not a no-go on the
    value three.
    First, the carrier selection: no result of the programme derives
    $V_{\rho} \cong \mathbb{C}^{2}$ from Born--Infeld parity; the parity input of
    O18~\cite{Beau2026a22} is a covariance of the response family and does not force the
    fibre structure, so the spinor carrier is a supplied input.
    One candidate source for that selection is now closed exactly.
    The fingerprint vectors of the O12 shell-span filtration are pure Fourier modes, so each
    proper shell span is a coordinate subspace whose Weyl support is a single line; its exact
    stabiliser in the Weil image is $U \rtimes M_{n}$, contained in the corresponding Borel
    subgroup.
    Since no binary polyhedral group embeds in a metacyclic Borel, that filtration selects
    neither a two-dimensional spinor carrier nor any of $2T$, $2O$, $2I$, for every odd prime,
    every generic block and every proper nonzero level.
    This closes one source of the carrier, not the carrier itself, and is not a no-go on the
    value three.
    Second, the observable identification: no theorem identifies the cumulative
    Gram--Schmidt span $\Sigma_{c}$ of O21 with the dimension of the neutral traceless
    module; since $n_{3}$ is defined as the shell where $\Sigma_{c}$ reaches three, the
    threshold is a supplied selection rule, not a derived quantity.
    A structural caution accompanies the positive result: a symmetric neutral generating
    set need not span three axes.
    In the faithful two-dimensional representation of $Q_{8}$, the set
    $\{\pm\mathbf{i}, \pm\mathbf{j}\}$ generates the group with only two generator axes;
    the third direction arises only under commutator closure, an operation not supplied by
    admissibility.
    The ADE spectral observation of three non-trivial eigenvalue classes is retained as a
    consistency check compatible with, but not probative of, the threefold structure.
  \end{abstract}

  \paragraph{Keywords.}
    spectral admissibility; spinor carrier; quaternions; Born–Infeld constraint;
    SU(2) representations; imaginary quaternions; Weil representation; ADE classification;
    Borel subgroup; binary polyhedral groups; carrier selection

    \clearpage
    \tableofcontents

    \clearpage
  \section{Position of the Problem after O21--O22}
  \label{sec:position}

  \subsection{What O21 established and what it left open}
    \label{ssec:o21-status}

    O21~\cite{Beau2026a25} established the shell-level saturation mechanism and formulated
    a conditional fibre-to-rate prescription involving $\dpair$ and $\bstar$.
    The structural analysis of the present paper is independent of that prescription, which
    has no native Heisenberg growth carrier~\cite{Beau2026sgn}.
    A central ingredient was the threshold condition
    \begin{equation}
      \label{eq:threshold}
      \Sigma_{c}(n_{3}) = 3,
    \end{equation}
    which defines the observable shell $n_{3}$ used throughout the cascade analysis.
    O21 used~\eqref{eq:threshold} as structural input: the value $3$ was interpreted as
    the number of physically stable directions selected by $\Pi$ in the Weil realisation
    of $\Heis$, but was not derived from the Born--Infeld structure.
    O21~§6.1 explicitly identified the derivation of~\eqref{eq:threshold} as the open
    problem of O23.

    O22~\cite{Beau2026a26} resolved the complementary question.
    Theorem~4.8 of O22 shows that projection locking under Born--Infeld spectral
    admissibility forces saturation to occur on a BFS shell ($n_{\mathrm{sat}} \in
  \mathbb{N}$), answering ``why a shell?''
    The present paper determines what the programme can and cannot currently say about
    ``why three?'': it proves the threefold dimension of the neutral sector of a supplied
    spinor carrier, and it states the two bridges that separate this theorem from a
    derivation of~\eqref{eq:threshold}.

  \subsection{The three distinct ``3''s that must not be conflated}
    \label{ssec:three-threes}

    The number $3$ appears in three structurally distinct places in the programme, and
    O23 must keep them separate throughout.

    \begin{enumerate}
      \item $\dim_{\mathbb{R}} \ImH = 3$: the real dimension of the imaginary quaternions,
      i.e.\ the subspace of $\mathbb{H}$ orthogonal to $\mathbb{R} \cdot \mathbf{1}$ and
      closed under the commutator.

      \item $\dim_{\mathbb{R}} \su(2) = 3$: the dimension of the Lie algebra $\su(2)$
      (equivalently, the dimension of $\SU(2)$ as a Lie group).
      This coincides with~(1) because $\su(2) \cong \ImH$ as real Lie algebras.

      \item $\Sigma_{c}(n_{3}) = 3$: the observable threshold of O21.
    \end{enumerate}

    The status of the links is the central point of this paper.
    The identification $(1) \cong (2)$ is a theorem, proved for a supplied irreducible
    two-dimensional carrier (Theorem~\ref{thm:adjoint-dimension}).
    The link $(2) \Rightarrow (3)$ is an open bridge: no map, intertwiner, or rank theorem
    identifies the Gram--Schmidt capacity count $\Sigma_{c}$ with the dimension of the
    neutral traceless module (Section~\ref{ssec:sigma-bridge}).
    Collapsing the three occurrences into a single unexplained ``$3$'' would reproduce
    exactly the conflation this paper is designed to prevent.

  \subsection{Scope and what this paper does not claim}
    \label{ssec:scope}

    O23 proves that a spinorial neutral sector carries exactly three independent
    directions: the traceless neutral sector of an irreducible two-dimensional
    $\SU(2)$-valued carrier is $\su(2) \cong \ImH$, of real dimension three.
    O23 does \emph{not} derive the carrier $V_{\rho} \cong \mathbb{C}^{2}$ from
    Born--Infeld parity, does \emph{not} identify $\Sigma_{c}$ with the dimension of the
    neutral module, and therefore does \emph{not} derive the threshold
    condition~\eqref{eq:threshold}.
    The value three in~\eqref{eq:threshold} has the status of a supplied selection rule,
    compatible with the spinorial structure, awaiting the two bridges of
    Section~\ref{sec:bridges}.
    It also does not determine the numerical value of $n_{3}$ as a function of
    $(c_{\mathrm{BI}}, \lambda_{n}, \text{graph geometry})$; that determination requires
    the extended numerical programme described in O22~§6.2 and remains an open direction.

  \section{The Carrier Setting and Its Support}
  \label{sec:exclusion}

  \subsection{Setup: neutral generators and the spinorial constraint}
    \label{ssec:neutral-setup}

    We recall the setup of the Admissibility Sub-programme~\cite{Beau2026a1,Beau2026a2}.
    Let $G$ be a finite group, $S = S^{-1} \subset G$ a symmetric generating set, and
    $\rho \colon G \to \SU(2)$ an irreducible two-dimensional representation.
    The \emph{neutrality condition} is
    \begin{equation}
      \label{eq:neutrality}
      \chi_{\rho}(s) := \mathrm{Tr}\,\rho(s) = 0 \quad \forall\, s \in S.
    \end{equation}
    This condition is the spectral admissibility constraint under bounded relational flux:
    it selects the generators that remain compatible with the Born--Infeld capacity
    bound~\cite{Beau2026a1,Beau2026a2}.
    By the Pauli tracelessness identity, each $\rho(s)$ satisfying~\eqref{eq:neutrality}
    can be written
    \begin{equation}
      \label{eq:pauli-axis}
      \rho(s) = \mathrm{i}(\vec{u}_{s} \cdot \vec{\sigma}), \quad
      \vec{u}_{s} \in S^{2} \subset \mathbb{R}^{3},
    \end{equation}
    where $\vec{\sigma} = (\sigma_{x}, \sigma_{y}, \sigma_{z})$ are the Pauli matrices.
    Each neutral generator therefore determines a unique axis $\pm \vec{u}_{s} \in
  \mathbb{RP}^{2}$~\cite{Beau2026a3}.

    The irreducible two-dimensional carrier is a \emph{supplied input} of this setting.
    The remainder of this section records what the assumption implies and, equally
    importantly, what it is not implied by.

  \subsection{Non-abelian support of the carrier}
    \label{ssec:abelian-exclusion}

    \begin{lemma}[Non-abelian support]
      \label{lem:abelian-exclusion}
      If a finite group $G$ admits an irreducible complex representation of dimension two,
      then $G$ is non-abelian.
    \end{lemma}

    \begin{proof}
      Every irreducible complex representation of a finite abelian group is
      one-dimensional, by Schur's lemma over $\mathbb{C}$.
    \end{proof}

    \begin{remark}[Neutrality alone does not exclude abelian groups]
      \label{rem:neutrality-not-enough}
      The non-abelian support is carried entirely by the irreducibility assumption, not by
      the neutrality condition.
      Without irreducibility, abelian groups do admit faithful two-dimensional
      representations with neutral symmetric generating sets: for $G = \mathbb{Z}/4$ with
      generator $s$, the representation
      $\rho = \chi_{1} \oplus \chi_{2}$ with $\chi_{1}(s) = \mathrm{i}$ and
      $\chi_{2}(s) = -\mathrm{i}$ is faithful, and
      $\chi_{\rho}(s) = \chi_{\rho}(s^{-1}) = 0$ on the symmetric generating set
      $S = \{s, s^{-1}\}$.
      Any exclusion of abelian realisations must therefore invoke the irreducibility of
      the carrier, which is part of the supplied setting rather than a consequence of
      admissibility.
    \end{remark}

    \begin{remark}[Candidate groups]
      \label{rem:exclusion-scope}
      Given the $\SU(2)$-valued irreducible two-dimensional carrier fixed in
      Section~\ref{ssec:neutral-setup}, the candidate groups are the finite non-abelian
      subgroups of $\SU(2)$ admitting a faithful irreducible two-dimensional
      representation.
      By the Klein--Schur classification recalled in~\cite{Beau2026a3}, these are the
      dicyclic groups $\mathrm{Dic}_{n}$ ($n \geq 2$), the binary tetrahedral group $2T$,
      the binary octahedral group $2O$, and the binary icosahedral group $2I$ (the cyclic
      subgroups of $\SU(2)$ are abelian and excluded by Lemma~\ref{lem:abelian-exclusion};
      outside the $\SU(2)$ setting, further unitary carriers exist and are not classified
      here).
      The binary tetrahedral group $2T$ is further structurally excluded by
      SpectralGramRigidity~\cite{Beau2026a3} (its full neutral sector generates
      only $Q_{8}$, not $2T$).
      The admissible candidates are therefore $Q_{8}$,
      $\mathrm{Dic}_{n}$ ($n \geq 3$), $2O$, and $2I$.
    \end{remark}

  \subsection{Associativity of the Weil framework and its limits}
    \label{ssec:octonion-exclusion}

    \begin{lemma}[Octonion incompatibility with the Weil realisation]
      \label{lem:octonion-exclusion}
      The octonion algebra $\mathbb{O}$ is incompatible with the class of unitary
      representations used by the O-series: the Weil representation on $\Heis$.
    \end{lemma}

    \begin{proof}
      The Weil representation $\rho_{c}$ on $\Heis$ is defined via the cocycle rule
      \[
        \rho_{c}(g_{1} g_{2}) = \phi(g_{1}, g_{2})\, \rho_{c}(g_{1})\, \rho_{c}(g_{2}),
      \]
      where $\phi$ is a $\mathrm{U}(1)$-valued two-cocycle.
      Since $\rho_{c}$ acts on a finite-dimensional complex Hilbert space $V$ via linear
      operators, its representation matrices are elements of $\mathrm{End}(V)$, which is
      an associative algebra by definition of operator composition.
      The cocycle identity forces
      $(\rho_{c}(g_{1})\, \rho_{c}(g_{2}))\, \rho_{c}(g_{3})
      = \rho_{c}(g_{1})\,(\rho_{c}(g_{2})\, \rho_{c}(g_{3}))$.
      The octonion algebra $\mathbb{O}$ is non-associative: there exist
      $x, y, z \in \mathbb{O}$ with $(xy)z \neq x(yz)$.
      A representation taking values in $\mathbb{O}$ cannot satisfy the cocycle identity
      for all triples $(g_{1}, g_{2}, g_{3})$.
      Hence no Weil-type representation on $\Heis$ can be realized over $\mathbb{O}$.
    \end{proof}

    \begin{remark}[Associativity does not select $\mathbb{H}$]
      \label{rem:no-hurwitz-selection}
      Lemma~\ref{lem:octonion-exclusion} excludes octonionic realisations; it does not
      select the quaternions.
      The representation matrices live in $\mathrm{End}(V) = M_{q}(\mathbb{C})$, an
      associative algebra that is \emph{not} a division algebra, so the Hurwitz
      classification of real normed division algebras~\cite{Hurwitz1898} does not apply
      to it and cannot force the representation algebra to be $\mathbb{H}$.
      The quaternionic structure of this paper enters only through the supplied
      $\SU(2)$ carrier: finite subgroups of $\SU(2)$ embed in the unit quaternions, and
      the traceless sector of the carrier is $\su(2) \cong \ImH$
      (Theorem~\ref{thm:adjoint-dimension}).
      Where the Hurwitz classification is used below, it is used only to position
      $\mathbb{H}$ among the alternatives \emph{given} an associative division structure,
      never to derive that structure from the Weil framework.
    \end{remark}

  \section{Three Neutral Directions from the Spinor Carrier}
  \label{sec:quaternionic}

  \subsection{The adjoint-dimension theorem}
    \label{ssec:main-theorem}

    \begin{theorem}[Three neutral directions from a spinor carrier]
      \label{thm:adjoint-dimension}
      Let $\rho \colon G \to \SU(2)$ be an irreducible two-dimensional representation of
      a finite group $G$ on $V_{\rho} \cong \mathbb{C}^{2}$, as fixed in
      Section~\ref{ssec:neutral-setup}.
      Then the traceless anti-Hermitian sector of the carrier,
      \[
        \su(V_{\rho}) := \{A \in \mathrm{End}(V_{\rho}) : A^{\dagger} = -A,\
        \mathrm{Tr}\,A = 0\},
      \]
      is a real Lie algebra isomorphic to $\su(2) \cong \ImH$, and
      \[
        \dim_{\mathbb{R}} \su(V_{\rho}) = 3.
      \]
      Every neutral generator image $\rho(s) = \mathrm{i}(\vec{u}_{s} \cdot \vec{\sigma})$
      lies on the unit sphere of this three-dimensional space, for the inner product
      $\langle A, B \rangle = \tfrac12 \mathrm{Tr}(A^{\dagger} B)$ that makes
      $\{\mathrm{i}\sigma_{x}, \mathrm{i}\sigma_{y}, \mathrm{i}\sigma_{z}\}$ orthonormal.
    \end{theorem}

    \begin{proof}
      Fix a unitary basis of $V_{\rho}$, so $\mathrm{End}(V_{\rho}) \cong
      M_{2}(\mathbb{C})$.
      The anti-Hermitian traceless matrices form the real Lie algebra $\su(2)$, with real
      basis $\{\mathrm{i}\sigma_{x}, \mathrm{i}\sigma_{y}, \mathrm{i}\sigma_{z}\}$, hence
      real dimension three.
      The map $\mathbf{i} \mapsto -\mathrm{i}\sigma_{x}$, $\mathbf{j} \mapsto
      -\mathrm{i}\sigma_{y}$, $\mathbf{k} \mapsto -\mathrm{i}\sigma_{z}$ extends to a real
      Lie algebra isomorphism $\ImH \to \su(2)$: the quaternion relations
      $[\mathbf{i},\mathbf{j}] = 2\mathbf{k}$ (cyclically) are matched exactly by
      $[-\mathrm{i}\sigma_{x}, -\mathrm{i}\sigma_{y}] = -[\sigma_{x},\sigma_{y}] =
      -2\mathrm{i}\sigma_{z} = 2(-\mathrm{i}\sigma_{z})$, with no rescaling of the
      bracket, realising the classical identification $\su(2) \cong \mathfrak{sp}(1)
      \cong \ImH$.
      Membership of the images in $\SU(2)$ makes~\eqref{eq:pauli-axis} applicable: each
      neutral generator image is $\mathrm{i}(\vec{u}_{s} \cdot \vec{\sigma})$ with
      $|\vec{u}_{s}| = 1$, an element of norm one for the stated inner product.
    \end{proof}

    \begin{remark}[What carries the ``3'']
      \label{rem:what-carries-three}
      Theorem~\ref{thm:adjoint-dimension} is a statement about the supplied carrier.
      The integer three is the real dimension of the adjoint (traceless) module of
      $\mathbb{C}^{2}$; equivalently $\su(2) \cong \operatorname{Sym}^{2}(V_{\rho})$ as
      real/adjoint modules.
      Nothing in the theorem derives the carrier itself: the selection of
      $V_{\rho} \cong \mathbb{C}^{2}$ from Born--Infeld admissibility is
      Bridge~1 of Section~\ref{sec:bridges} and is open.
    \end{remark}

    \begin{remark}[The ``$3$'' is not the dimension of $\mathbb{H}$]
      \label{rem:not-four}
      The integer $3$ in Theorem~\ref{thm:adjoint-dimension} is $\dim_{\mathbb{R}}
      \ImH$, not $\dim_{\mathbb{R}} \mathbb{H} = 4$.
      The full quaternion algebra contributes one additional real dimension ($\mathbb{R}
      \cdot \mathbf{1}$), which is central and plays no role in the neutral traceless
      sector.
      The neutrality condition~\eqref{eq:neutrality} projects out this
      commutative direction, leaving the three-dimensional traceless sector.
    \end{remark}

  \subsection{Generator axes need not span the three directions}
    \label{ssec:generator-axes}

    The identification of ``stable admissible directions'' with ``dimensions of the
    neutral sector'' requires care: the axes realised by a given neutral generating set
    can span a proper subspace of $\ImH$.

    \begin{proposition}[Two-axis neutral generating set of $Q_{8}$]
      \label{prop:q8-two-axes}
      In the faithful two-dimensional representation of $Q_{8} \subset \SU(2)$, the
      symmetric set $S = \{\pm\mathbf{i}, \pm\mathbf{j}\}$ generates $Q_{8}$, consists of
      traceless (neutral) elements, and its generator axes span only the two-dimensional
      plane $\operatorname{span}_{\mathbb{R}}\{\mathbf{i}, \mathbf{j}\} \subset \ImH$.
    \end{proposition}

    \begin{proof}
      The products of $\pm\mathbf{i}$ and $\pm\mathbf{j}$ include
      $\mathbf{i}\mathbf{j} = \mathbf{k}$ and $\mathbf{i}^{2} = -1$, so
      $\langle S \rangle = \{\pm 1, \pm\mathbf{i}, \pm\mathbf{j}, \pm\mathbf{k}\} =
      Q_{8}$, and $S$ is symmetric since $\mathbf{i}^{-1} = -\mathbf{i}$,
      $\mathbf{j}^{-1} = -\mathbf{j}$.
      In the standard representation, $\mathbf{i}, \mathbf{j}$ map to
      $\mathrm{i}\sigma_{x}, \mathrm{i}\sigma_{y}$ (up to labelling), which are traceless.
      The axes of the four generators are $\pm\vec{e}_{x}$ and $\pm\vec{e}_{y}$, spanning
      the $\{\mathbf{i},\mathbf{j}\}$-plane only.
    \end{proof}

    \begin{remark}[Closure is an extra operation]
      \label{rem:closure-extra}
      The third direction $\mathbf{k}$ is generated by the \emph{commutator} of the two
      listed axes, i.e.\ by the Lie closure of the generator set, not by the generator
      axes themselves.
      Proposition~\ref{prop:q8-two-axes} therefore rules out the inference from ``$S$ is
      an admissible neutral generating set'' to ``$S$ realises three independent axes'':
      a count of generator axes can stop at two.
      Whether the capacity observable $\Sigma_{c}$ performs the Lie closure --- i.e.\
      whether the Gram--Schmidt span at saturation measures the closed algebra rather
      than the listed axes --- is precisely Bridge~2 of Section~\ref{sec:bridges} and is
      not established.
      Independence of axes also does not by itself give the quaternion relations: for
      matrices $\mathrm{i}(\vec{u} \cdot \vec{\sigma})$ and
      $\mathrm{i}(\vec{v} \cdot \vec{\sigma})$, anticommutation requires
      $\vec{u} \cdot \vec{v} = 0$, not merely $\vec{u} \nparallel \vec{v}$.
    \end{remark}

  \section{The Two Open Bridges}
  \label{sec:bridges}

  \subsection{Bridge 1: carrier selection}
    \label{ssec:carrier-bridge}

    Theorem~\ref{thm:adjoint-dimension} presupposes the irreducible two-dimensional
    carrier.
    No result of the programme currently derives this carrier from the Born--Infeld
    structure.
    O18~\cite{Beau2026a22} establishes that Born--Infeld parity acts on the effective
    responses as a covariance --- even-order responses preserved under the transported-probe
    comparison, odd-order responses reversed (O18 Proposition~2.6) --- and that evenness of
    the action does not force parity-related configurations into a common fibre
    (O18 Proposition~2.4);
    the fibre identification is the typed classification problem O18 Problem~2.8, whose
    data are not supplied.
    Consequently, the chain ``Born--Infeld parity $\Rightarrow$ spinor carrier'' is not
    available: the carrier $V_{\rho} \cong \mathbb{C}^{2}$ is a supplied input, motivated
    by the Weil realisation and by the Gram rigidity results of~\cite{Beau2026a3}, and any
    downstream use must carry this status.

  \subsection{Bridge 2: observable identification}
    \label{ssec:sigma-bridge}

    O21~\cite{Beau2026a25} defines $\Sigma_{c}$ as a cumulative Gram--Schmidt span
    dimension in the fingerprint pipeline, and defines $n_{3}$ as the shell at which
    $\Sigma_{c}$ reaches the value three.
    No theorem identifies $\Sigma_{c}$ at saturation with
    $\dim_{\mathbb{R}} \su(V_{\rho})$: the programme supplies no map, intertwiner, rank
    identity, or spectral theorem connecting the Gram--Schmidt capacity count to the
    dimension of the neutral traceless module.
    Since $n_{3}$ is defined by the threshold value three, the statement
    ``$\Sigma_{c}(n_{3}) = 3$'' is definitional at the level of O21 and cannot, by
    itself, certify that the projection selects a three-dimensional module.
    The correct status of the threshold is therefore: \emph{supplied selection rule,
    compatible with the spinorial structure of Theorem~\ref{thm:adjoint-dimension},
    awaiting an observable-identification theorem}.

  \subsection{The honest chain}
    \label{ssec:honest-chain}

    The status-typed chain replacing the former derivation claim is:
    \begin{align}
      \label{eq:chain}
      &\text{BI parity (covariant, proved)}
      \;\xrightarrow{\ \text{open: O18 Problem 2.8}\ }\;
      \text{pair fibre } c \leftrightarrow q-c \notag\\
      &\;\xrightarrow{\ \text{open: Bridge 1, not from the O12 filtration (Cor.~\ref{cor:no-binary-polyhedral})}\ }\;
      V_{\rho} \cong \mathbb{C}^{2}
      \;\xrightarrow{\ \text{proved: Thm~\ref{thm:adjoint-dimension}}\ }\;
      \dim_{\mathbb{R}} \su(V_{\rho}) = 3
      \;\xrightarrow{\ \text{open: Bridge 2}\ }\;
      \Sigma_{c}(n_{3}) = 3.
    \end{align}
    Only the third arrow is a theorem.
    Subsection~\ref{ssec:o12-obstruction} proves that one candidate source for the second
    arrow, the deposited O12 shell-span filtration, cannot supply it.

  \subsection{A proved obstruction on Bridge 1: the O12 filtration cannot supply the carrier}
    \label{ssec:o12-obstruction}

    Bridge~1 asks where a non-circular selection of $V_{\rho} \cong \mathbb{C}^{2}$ could come
    from.
    A natural candidate is the deposited O12 shell-span filtration itself: if some subgroup of
    the Weil image stabilised it, that subgroup would be selected by the construction rather
    than supplied by hand, and a binary polyhedral stabiliser would carry a two-dimensional
    spinor module.
    This subsection proves that this candidate fails, for a structural reason that no choice of
    prime, block or depth can evade.

    Fix an odd prime $q$ and a generic block $c = (c_{1}, c_{2}, c_{3})$ in the sense of O12,
    that is $c_{i} \not\equiv 0$ and $\sum_{i} c_{i} \not\equiv 0 \pmod{q}$.
    Let $W_{<n} \subset \mathbb{C}^{q}$ denote the span of the $k = 3$ fingerprint vectors of
    the BFS shells of depth less than $n$, and let $P_{n}$ be the orthogonal projector onto
    $W_{<n}$.
    The Weil representation $\rho_{c}$ of $\SL(2,\mathbb{Z}/q\mathbb{Z})$ acts on the same
    $\mathbb{C}^{q}$ and normalises the Heisenberg action~\cite{Weil1964,Gerardin1977}.
    Write $W_{c}(v)$, $v \in \mathbb{F}_{q}^{2}$, for the symmetrised Weyl operators, so that
    $\rho_{c}(g) W_{c}(v) \rho_{c}(g)^{-1} = W_{c}(gv)$.

    \begin{lemma}[Fingerprint vectors are pure Fourier modes]
      \label{lem:pure-fourier}
      Every $k = 3$ fingerprint vector of the O12 construction is a scalar multiple of a single
      Fourier mode $e_{f}$, with $f = \sum_{i} c_{i} b_{i} \bmod q$, where $b_{i}$ is the
      $b$-component of the $i$-th partial Heisenberg product.
      Consequently $W_{<n} = \mathrm{span}\{ e_{f} : f \in F_{n} \}$ is a coordinate subspace in
      the Fourier basis, where $F_{n} \subset \mathbb{F}_{q}$ is the exact set of frequencies
      reached at depth less than $n$, and $\dim W_{<n} = |F_{n}|$.
    \end{lemma}

    \begin{proof}
      A fingerprint vector is the pointwise product of three Schr\"{o}dinger-model vectors, the
      $i$-th carrying the phase $c_{i}(\gamma_{i} + b_{i}(x - a_{i}))$.
      Summing the three phases, the dependence on $x$ is $\big(\sum_{i} c_{i} b_{i}\big) x$ and
      every remaining term is independent of $x$.
      The product is therefore a constant phase times $e_{f}$ with $f = \sum_{i} c_{i} b_{i}$.
      A span of Fourier modes is a coordinate subspace in that basis, and its dimension is the
      number of distinct modes it contains.
    \end{proof}

    Let $L = \{(a,0) : a \in \mathbb{F}_{q}\}$, let $B(L)$ be its stabiliser in
    $\SL(2,\mathbb{Z}/q\mathbb{Z})$, a Borel subgroup of order $q(q-1)$, let
    $U = \{ \left(\begin{smallmatrix} 1 & t \\ 0 & 1 \end{smallmatrix}\right) \}$ be its
    unipotent radical, and set
    \[
      M_{n} = \{ s \in \mathbb{F}_{q}^{\times} : s F_{n} = F_{n} \}.
    \]

    \begin{theorem}[Exact stabiliser of the O12 filtration]
      \label{thm:o12-stabiliser}
      For every proper nonzero level, that is whenever $0 < |F_{n}| < q$,
      \[
        \mathrm{Stab}(W_{<n})
        = \{ g \in \SL(2,\mathbb{Z}/q\mathbb{Z}) : \rho_{c}(g) W_{<n} = W_{<n} \}
        = U \rtimes M_{n}
        \subseteq B(L).
      \]
    \end{theorem}

    \begin{proof}
      Expand $P_{n} = q^{-1} \sum_{v} p_{n}(v) W_{c}(v)$ in the Weyl operator basis of
      $\mathrm{End}(\mathbb{C}^{q})$, with $p_{n}(v) = \mathrm{tr}(P_{n} W_{c}(v)^{\dagger})$.
      Since $\rho_{c}(g)$ is unitary, $\rho_{c}(g) W_{<n} = W_{<n}$ is equivalent to
      $\rho_{c}(g) P_{n} \rho_{c}(g)^{\dagger} = P_{n}$, and by the intertwining relation this
      holds if and only if $p_{n} \circ g^{-1} = p_{n}$ on $\mathbb{F}_{q}^{2}$.
      By Lemma~\ref{lem:pure-fourier}, $P_{n}$ is diagonal in the Fourier basis, so $p_{n}$ is
      supported on $L$.
      Any $g$ preserving $p_{n}$ preserves its support, hence preserves $L$, so
      $\mathrm{Stab}(W_{<n}) \subseteq B(L)$; the hypothesis $|F_{n}| < q$ is what makes $p_{n}$
      non-constant off the origin, and $|F_{n}| > 0$ excludes the zero projector.
      Within $B(L)$, write $g$ as a product of a unipotent and a torus element.
      The unipotent $U$ is diagonal in the Fourier basis, hence preserves every coordinate
      subspace and lies in the stabiliser unconditionally.
      The torus element $\mathrm{diag}(s, s^{-1})$ sends $e_{f}$ to $e_{sf}$, hence preserves
      $W_{<n}$ if and only if $s F_{n} = F_{n}$, that is $s \in M_{n}$.
    \end{proof}

    \begin{corollary}[No binary polyhedral carrier is selected]
      \label{cor:no-binary-polyhedral}
      For every odd prime $q$, every generic block and every proper nonzero level, the
      stabiliser of $W_{<n}$ contains no subgroup isomorphic to $2T \cong \SL(2,3)$,
      $2O$ or $2I \cong \SL(2,5)$, and is not itself isomorphic to any of them.
      In particular the O12 shell-span filtration does not select an irreducible
      two-dimensional carrier $V_{\rho} \cong \mathbb{C}^{2}$, and the choice of such a carrier
      by restriction to a binary polyhedral subgroup remains supplied structure.
    \end{corollary}

    \begin{proof}
      By Theorem~\ref{thm:o12-stabiliser} the stabiliser lies in $B(L)$, which is metacyclic:
      it is an extension of the cyclic torus of order $q-1$ by the cyclic unipotent radical of
      order $q$.
      Every subgroup of a metacyclic group is metacyclic.
      The binary polyhedral groups $2T$, $2O$ and $2I$ are not metacyclic, since each has a
      non-cyclic derived subgroup, so none embeds in $B(L)$.
      By Dickson's classification of the subgroups of $\SL(2,q)$~\cite{Dickson1901}, these
      groups do occur in $\SL(2,\mathbb{Z}/q\mathbb{Z})$ for suitable $q$ --- $2I$ whenever
      $q \equiv \pm 1 \pmod{5}$ --- so the obstruction is not an availability accident:
      the groups are present in the ambient group and are simply not the ones the filtration
      selects.
    \end{proof}

    \begin{remark}[What is universal, what is measured, and what is not established]
      \label{rem:o12-status}
      Three statements of different epistemic strength must be kept apart.
      Theorem~\ref{thm:o12-stabiliser} and Corollary~\ref{cor:no-binary-polyhedral} are
      universal: they hold for every odd prime, every generic block and every proper nonzero
      level, and they carry the whole negative conclusion.
      The \emph{value} of $M_{n}$ is not universal.
      On the documented sampled blocks at $q = 53, 101, 211$ the exact multiplier groups are
      $M_{n} = \{\pm 1\}$ at every proper level, so that $\mathrm{Stab}(W_{<n}) \cong C_{2q}$
      with normaliser $B(L)$; but genericity does not force this, and generic blocks with
      $|M_{n}| \in \{4, 6\}$ at some depth exist, the excess being transient in depth.
      Over $2000$ further generic blocks per prime and two independent sampling streams, the
      proportion of blocks exhibiting $M_{n} \neq \{\pm 1\}$ at some level of depth at most six
      was $25/2000$, $7/2000$, $1/2000$ and $26/2000$, $4/2000$, $0/2000$ at
      $q = 53, 101, 211$ respectively.
      That the intersection over levels is $\{\pm 1\}$ for \emph{all} generic blocks is
      \emph{not established}: it held in every one of the $12\,008$ draws performed, and no
      proof is offered.
      The negative conclusion of Corollary~\ref{cor:no-binary-polyhedral} does not depend on
      any of these measurements.
    \end{remark}

    \begin{remark}[Scope of the obstruction]
      \label{rem:o12-scope}
      Corollary~\ref{cor:no-binary-polyhedral} rules out one \emph{source} for the carrier, not
      the carrier itself, and it is not a no-go on the value three.
      Bridge~1 remains open: nothing here proves that no admissible construction selects a
      spinor carrier, only that this filtration does not.
      Two consequences are worth recording for future attempts.
      First, enriching the object cannot help: if $g$ fixes a set of vectors, a labelled shell
      family, or a configuration carrying the BFS metric, then it fixes their linear span, so
      the stabiliser of any such configuration is contained in $\mathrm{Stab}(W_{<n})$ and
      therefore in $B(L)$.
      Second, the obstruction is exactly the one-dimensionality of the Weyl support, so a
      replacement construction must produce a proper nonzero projector whose Weyl support spans
      at least two distinct lines through the origin.
      That condition is necessary, not sufficient.
    \end{remark}

    \paragraph{Reproduction.}
      The statements of this subsection are proved above and are independently verified by the
      diagnostic in the \texttt{code/} directory of this paper's repository, which computes the
      exact frequency sets $F_{n}$, the multiplier groups $M_{n}$ and the stabilisers by integer
      arithmetic, and cross-checks them against an independent floating-point computation of the
      projectors, aborting on any disagreement.
      The exact commands, thresholds and seeds are documented in the repository README.

  \section{Spectral Observations}
  \label{sec:spectral}

  \subsection{$Q_{8}$ as the minimal isotropic prototype}
    \label{ssec:q8}

    The quaternion group $Q_{8}$ is the smallest finite subgroup of $\SU(2)$ admitting a
    neutral generating set, and SpectralGramRigidity~\cite{Beau2026a3} identifies
    it as the unique admissible group with an orthogonal Gram pattern
    ($G_{st} = 0$ for all $t \neq s^{\pm 1}$).
    Its six traceless elements $\{\pm\mathbf{i}, \pm\mathbf{j}, \pm\mathbf{k}\}$
    pair into three conjugate pairs, each determining one of the three coordinate axes
    of $\ImH$, with the second-moment tensor $M \propto I$~\cite{Beau2026a3}.
    The \emph{full} neutral sector of $Q_{8}$ thus realises the three directions of
    Theorem~\ref{thm:adjoint-dimension} isotropically; its two-axis generating subset
    (Proposition~\ref{prop:q8-two-axes}) shows that a generating set alone need not.

  \subsection{Three-level ADE stratigraphic structure (consistency observation)}
    \label{ssec:ade-three-levels}

    For the ADE binary Cayley graphs physically relevant to the Cosmochrony
    programme --- in particular $2I$ with order-$4$ and order-$5$ generators --- the
    Cayley eigenvalue formula gives, for each irreducible representation $\rho$ of $G$,
    \[
      \lambda_{\rho} = |S| - \frac{1}{d_{\rho}} \sum_{s \in S} \chi_{\rho}(s).
    \]
    The character sum $\sum_{s \in S} \chi_{\rho}(s) / d_{\rho}$ takes exactly three
    distinct values over the non-trivial irreducible representations of $G$, corresponding
    to three distinct eigenvalue classes $\lambda_{1} < \lambda_{2} < \lambda_{3}$, and
    hence three stratigraphic levels~\cite{Beau2026a4}.

    This observation is a consistency check: it is compatible with a threefold neutral
    structure, and a derived threefold structure would explain it.
    It is not probative of the threshold~\eqref{eq:threshold}: the count of distinct
    character-sum values is a property of the chosen groups and generating sets, and no
    theorem connects it to the Gram--Schmidt capacity observable $\Sigma_{c}$.

  \section{Discussion}
  \label{sec:discussion}

  \subsection{What O23 establishes and what it does not}
    \label{ssec:scope-discussion}

    O23 establishes, unconditionally for a supplied irreducible two-dimensional
    $\SU(2)$-valued carrier, that the neutral traceless sector is $\su(2) \cong \ImH$ with exactly three
    independent directions (Theorem~\ref{thm:adjoint-dimension}), and that generator-axis
    counting cannot replace this module statement
    (Proposition~\ref{prop:q8-two-axes}).
    It establishes further, and unconditionally, the exact stabiliser of the O12 shell-span
    filtration in the Weil image (Theorem~\ref{thm:o12-stabiliser}) and the resulting exclusion
    of every binary polyhedral group as a selected stabiliser
    (Corollary~\ref{cor:no-binary-polyhedral}).
    It does not derive the carrier, does not identify $\Sigma_{c}$ with the module
    dimension, and therefore does not derive $\Sigma_{c}(n_{3}) = 3$.
    The threshold retains the status of a supplied selection rule.

    A derivation of the threshold would have to close both bridges of
    Section~\ref{sec:bridges}: a non-circular selection of the spinor carrier from the
    projection structure, and an observable-identification theorem for $\Sigma_{c}$.
    Both are open problems of the programme, and the second requires in addition a rank
    measurement protocol that does not fix the target dimension in advance.
    What Corollary~\ref{cor:no-binary-polyhedral} contributes to the first is a proved
    restriction on where the carrier may come from, and a concrete invariant criterion --- a
    Weyl support spanning at least two lines --- that any replacement construction must satisfy
    before it is worth testing.

  \subsection{Relation to O21--O22}
    \label{ssec:conditional}

    O22's shell result (saturation occurs on a BFS shell) is unaffected.
    O21's use of $n_{3}$ remains well-defined as a definition; what changes is its
    epistemic label: the value three entering the definition is a supplied selection
    rule compatible with the spinorial structure, not a derived constant.
    Downstream consumers of ``$\Sigma_{c}(n_{3}) = 3$'' must cite it with this status.

  \section{Conclusion}
  \label{sec:conclusion}

  A spinorial neutral sector carries exactly three independent directions: for an
  irreducible two-dimensional $\SU(2)$-valued carrier $V_{\rho}$, the traceless neutral sector is
  $\su(2) \cong \ImH$, of real dimension three
  (Theorem~\ref{thm:adjoint-dimension}).
  This theorem explains why the value three is the natural threshold for a spinor-carried
  admissible sector, and the $Q_{8}$ two-axis generating set
  (Proposition~\ref{prop:q8-two-axes}) shows that the statement lives at the level of the
  neutral module, not of generator counting.

  The threshold condition $\Sigma_{c}(n_{3}) = 3$ of O21 is not derived by this theorem.
  Two bridges separate them, both open: the selection of the spinor carrier from the
  Born--Infeld projection structure, and the identification of the Gram--Schmidt capacity
  observable with the dimension of the neutral module.
  The threefold threshold is accordingly a supplied selection rule, compatible with and
  motivated by the spinorial structure, and the closure of either bridge is a
  well-defined open problem of the programme.

  On the first bridge, one candidate source is now excluded by an exact theorem.
  The O12 shell-span filtration has stabiliser $U \rtimes M_{n}$ inside a Borel subgroup of
  $\SL(2,\mathbb{Z}/q\mathbb{Z})$ for every proper level, every generic block and every odd
  prime (Theorem~\ref{thm:o12-stabiliser}), and no binary polyhedral group embeds in a
  metacyclic Borel, so that filtration selects no spinor carrier
  (Corollary~\ref{cor:no-binary-polyhedral}).
  The exclusion is sharpest where the groups are available: for $q \equiv \pm 1 \pmod 5$ the
  group $2I$ does sit inside $\SL(2,\mathbb{Z}/q\mathbb{Z})$, and the filtration still does not
  select it.

  \paragraph{Interpretive outlook (interpretation, not a result).}
    Read structurally, this says that the admissible filtration studied so far is too
    ``one-directional'' to carry a spin degree of freedom: its entire Weyl content lies along a
    single line of the phase-space plane, and a spinor needs at least two independent
    directions to exist.
    If the three neutral directions of Theorem~\ref{thm:adjoint-dimension} are to be more than
    a supplied convenience --- if they are eventually to explain why observable structure comes
    in threes --- then the object carrying them must be built from something with a genuinely
    two-dimensional support.
    This reading is offered as orientation for where to look next, not as an established
    consequence: the present paper proves where the carrier does \emph{not} come from, and says
    nothing about where it does.

  \bigskip
  \noindent
  \textbf{Acknowledgements.}
  This paper is part of the Spectral Admissibility Sub-programme of the Cosmochrony
  framework~\cite{Cosmochrony}.

  \appendix

  \section{Position of $\mathbb{H}$ among the Normed Division Algebras}
  \label{app:not-2-not-7}

  Given the supplied carrier setting, the position of the quaternions among the real
  normed division algebras $\{\mathbb{R}, \mathbb{C}, \mathbb{H}, \mathbb{O}\}$ of the
  Hurwitz classification~\cite{Hurwitz1898} can be summarised as follows.

  \paragraph{$\mathbb{C}$ (imaginary dimension 1).}
    A carrier valued in $\mathbb{C}^{\times}$ is one-dimensional over $\mathbb{C}$, hence
    corresponds to an abelian image and is outside the irreducible two-dimensional
    setting (Lemma~\ref{lem:abelian-exclusion}).
    Its imaginary subspace $\mathrm{Im}\,\mathbb{C} = \mathrm{i}\mathbb{R}$ offers a
    single neutral axis.

  \paragraph{$\mathbb{O}$ (imaginary dimension 7).}
    The octonions are non-associative, and Lemma~\ref{lem:octonion-exclusion} shows that
    the Weil cocycle structure cannot be realised over $\mathbb{O}$.
    The seven-dimensional imaginary space is therefore inaccessible within the Weil
    framework of the O-series.

  \paragraph{$\mathbb{H}$ (imaginary dimension 3).}
    The quaternions are the unique associative non-commutative entry, and the unit
    quaternions realise $\SU(2)$, the carrier group of the supplied setting; their
    imaginary subspace is the three-dimensional neutral sector of
    Theorem~\ref{thm:adjoint-dimension}.

  \paragraph{Caveat.}
    This positioning takes the associative division structure as given.
    It does not constitute a derivation: the Weil representation algebra
    $M_{q}(\mathbb{C})$ is not a division algebra, so the Hurwitz classification does not
    apply to it (Remark~\ref{rem:no-hurwitz-selection}), and the selection of the
    two-dimensional carrier is Bridge~1 of Section~\ref{sec:bridges}.

    \clearpage
    \bibliographystyle{plain}
    \bibliography{cosmochrony-bibliography,external-refs}

\end{document}
